%%
%% O5_airy_gap_problem.tex
%% Self-contained problem document for Open Node O5.
%% Subfield: Spectral theory of integrable operators /
%%           Random matrix edge universality / Airy operator analysis.
%% No prime arithmetic, no Gram matrices, no Kadec-type arguments.
%%
\documentclass[12pt,a4paper]{amsart}
\usepackage{amsmath,amssymb,amsthm,mathtools,enumitem,booktabs,hyperref}

\newtheorem{theorem}{Theorem}[section]
\newtheorem{lemma}[theorem]{Lemma}
\newtheorem{proposition}[theorem]{Proposition}
\newtheorem{corollary}[theorem]{Corollary}
\theoremstyle{definition}
\newtheorem{definition}[theorem]{Definition}
\newtheorem{remark}[theorem]{Remark}
\newtheorem{problem}[theorem]{Problem}
\newtheorem{conjecture}[theorem]{Conjecture}

\newcommand{\Ai}{\mathrm{Ai}}
\newcommand{\KAi}{K^{\mathrm{Airy}}}
\newcommand{\DAi}{D^{\mathrm{Airy}}}
\newcommand{\PAi}{\Pi_{\mathrm{Airy}}}
\newcommand{\ip}[2]{\langle #1,\,#2\rangle}
\newcommand{\norm}[1]{\|#1\|}

\title{O5: Spectral Gap of the Airy Mismatch Operator\\
{\normalsize A Self-Contained Problem Document}}
\author{Sebastian Schmalnauer}
\date{May 2026}

\begin{document}
\maketitle

\begin{abstract}
This document isolates Open Problem O5 of Paper~XXII as a
self-contained problem in spectral analysis of integrable operators.
No prime arithmetic, Gram matrices, frame theory, or
Riemann--Hilbert analysis is required.
The problem asks whether the Airy mismatch operator
$\DAi = \PAi(A^{\mathrm{arith}} - \KAi)\PAi$
has a spectral gap: $\inf\sigma(\DAi) \ge c_2 > 0$.
We formulate the precise statement, identify the relevant classical
tools (Airy oscillation, Fredholm theory, GUE edge universality),
and describe three independent attack routes.
\end{abstract}

\tableofcontents

%%=================================================================
\section{The Problem}
\label{sec:problem}
%%=================================================================

\subsection{The Airy kernel operator}

The \emph{Airy kernel} is:
\[
  \KAi(x,y)
  := \frac{\Ai(x)\Ai'(y) - \Ai'(x)\Ai(y)}{x-y},
  \quad x,y \in \mathbb{R}.
\]
The associated operator $K^{\mathrm{Airy},(s)}$
on $L^2([s,+\infty))$ is trace-class for every $s \in \mathbb{R}$
and its Fredholm determinant is the GUE edge gap probability:
$F_{\mathrm{TW}}(s) = \det(I - K^{\mathrm{Airy},(s)})$ (Tracy--Widom~\cite{TracyWidom1994}).

\subsection{The arithmetic sampling form in the Airy limit}

Let $\{p_j^*\}$ be the prime-scaled points (as in O4 document)
and let $A^{\mathrm{arith}}$ be the limiting arithmetic
sampling form in the Airy scaling limit
$x = s + c^{-2/3}\cdot(\text{edge displacement})$:
\[
  \ip{A^{\mathrm{arith}} f}{g}
  := \lim_{N\to\infty}\frac{1}{N}
  \sum_j f(p_j^{\mathrm{Ai}}) \overline{g(p_j^{\mathrm{Ai}})},
\]
where $p_j^{\mathrm{Ai}}$ are the primes mapped to the
Airy scale via the canonical edge scaling.

\begin{problem}[O5: Spectral gap of $\DAi$]\label{prob:O5}
Let $\PAi$ be the orthogonal projection onto the
closed linear span of $\{\Ai(\cdot - s) : s \in [0,R]\}$
(the Airy edge subspace on $[0,R]$ for some $R > 0$).
Define the \emph{Airy mismatch operator}:
\[
  \DAi
  := \PAi\,(A^{\mathrm{arith}} - \KAi)\,\PAi
  : \PAi L^2([0,R]) \to \PAi L^2([0,R]).
\]
Is $\DAi$ positive definite? More precisely:
\begin{equation}\label{eq:O5}
  \inf\,\sigma(\DAi) \ge c_2^{\mathrm{Airy}} > 0
\end{equation}
with $c_2^{\mathrm{Airy}}$ independent of $R$ and of
the specific prime sequence (depending only on
the Beurling density $D^+(\{p_j^{\mathrm{Ai}}\})$)?
\end{problem}

\begin{remark}[Structure of the problem]
$\DAi$ is the \emph{difference} between two objects:
\begin{itemize}
  \item $A^{\mathrm{arith}}$: discrete, arithmetic,
    governed by the prime distribution.
    Its spectrum has no universal structure.
  \item $\KAi$: continuous, universal,
    governed by Airy functions.
    Its eigenvalues are known (Tracy--Widom, GUE edge).
\end{itemize}
The positivity of their \emph{difference}
is a non-cancellation statement:
the arithmetic sampling operator must dominate
the Airy kernel operator in the Airy edge subspace.
\end{remark}

\begin{remark}[Consequence for Paper~XXII]
A positive answer implies, via the Airy scaling limit
convergence $\Dedge \to \DAi$ (Proposition~5.1 of Paper~XXII
conditional on Conjecture~3.3), the uniform edge coercivity
$\inf\sigma(\Dedge) \ge c_2^{\mathrm{Airy}}/2 > 0$.
This is the Main Theorem of Paper~XXII.
\end{remark}

%%=================================================================
\section{Key Structural Facts}
\label{sec:facts}
%%=================================================================

\begin{proposition}[Non-degeneracy of prime Airy sampling]%
\label{prop:nondeg}
The Airy function satisfies $\Ai(t) > 0$ for all $t > 0$.
In particular, for prime-scaled points $p_j^{\mathrm{Ai}} > 0$:
\[
  |\Ai(p_j^{\mathrm{Ai}})|^2 > 0
  \quad \text{for all } j.
\]
Therefore the diagonal of $A^{\mathrm{arith}}$
in the basis $\{\Ai(\cdot - a_k)\}$
(where $a_k < 0$ are Airy zeros) is strictly positive.
\end{proposition}

\begin{remark}[Why this is non-trivial]
Proposition~\ref{prop:nondeg} guarantees \emph{pointwise}
non-degeneracy. The problem asks for \emph{operator-level}
uniform positivity, which is a global statement about
all linear combinations $f = \sum_k \alpha_k \Ai(\cdot - a_k)$.
Pointwise non-degeneracy is necessary but not sufficient
for $\inf\sigma(\DAi) > 0$.
\end{remark}

\begin{proposition}[Spectrum of $\KAi$ is known]%
\label{prop:kairy_spectrum}
For $R < \infty$:
\begin{itemize}
  \item The eigenvalues of $K^{\mathrm{Airy},(0)}$
    on $L^2([0,R])$ lie in $[0,1)$ and accumulate only at $0$ and $1$.
  \item $\|K^{\mathrm{Airy},(0)}\|_{\mathfrak{S}_1}
    = \int_0^{\infty} K^{\mathrm{Airy},(0)}(x,x)\,dx < \infty$.
  \item $\det(I - \KAi^{(0)}) = F_{\mathrm{TW}}(0) \approx 0.9597$
    (numerically known to high precision~\cite{Bornemann2010}).
\end{itemize}
\end{proposition}

%%=================================================================
\section{Three Attack Routes}
\label{sec:attacks}
%%=================================================================

\subsection{Route A: Quantitative oscillation bound via Airy zeros}

\textbf{Goal.}
Bound $\ip{\DAi f}{f}$ from below for all unit vectors
$f \in \PAi L^2([0,R])$ using the oscillation structure
of Airy functions.

\textbf{Mechanism.}
Express $f = \sum_k \alpha_k \phi_k$ in the eigenbasis
$\{\phi_k\}$ of $\KAi$ (which is explicitly related
to Airy functions).
Use the asymptotics of Airy zeros:
$a_k \sim -(3\pi k/2)^{2/3}$ as $k \to \infty$.
Between consecutive Airy zeros,
$|\Ai(t)|^2$ has a single maximum;
the prime-scaled sampling points hit each such interval
(by the prime density argument),
providing a Gauss-quadrature-type lower bound on
$\ip{A^{\mathrm{arith}} f}{f}$.

\textbf{Key inequality to establish:}
\[
  \frac{1}{N}\sum_j |f(p_j^{\mathrm{Ai}})|^2
  \ge \int_0^R |f(t)|^2\,dt - \varepsilon(N)\norm{f}^2
\]
with $\varepsilon(N) \to 0$, plus
$\ip{\KAi f}{f} \le \mu_{\max}\norm{f}^2 < 1$,
which gives $\ip{\DAi f}{f} \ge (1 - \mu_{\max} - \varepsilon)\norm{f}^2 > 0$.

\subsection{Route B: Connection to the Tracy--Widom constant}

\textbf{Goal.}
Estimate $c_2^{\mathrm{Airy}}$ in terms of the
Tracy--Widom Fredholm determinant $F_{\mathrm{TW}}(s)$.

\textbf{Mechanism.}
The spectral gap of $\DAi$ is related to the
$\mathfrak{S}_1$-distance between $A^{\mathrm{arith}}$
and $\KAi$:
\[
  \inf\sigma(\DAi)
  \ge \lambda_{\min}(A^{\mathrm{arith}}|_{\mathrm{edge}})
    - \|\KAi|_{\mathrm{edge}}\|.
\]
The operator norm of $\KAi$ on $L^2([0,R])$ is bounded
by $1 - F_{\mathrm{TW}}(0) \approx 0.04$ (at $s=0$).
If $\lambda_{\min}(A^{\mathrm{arith}}|_{\mathrm{edge}}) > 0.04$,
then $c_2^{\mathrm{Airy}} > 0$.

\textbf{This is a quantitative route:}
it converts the abstract positivity question
into a numerically checkable bound
(cf.\ Bornemann~\cite{Bornemann2010} for high-precision
Tracy--Widom numerics).

\subsection{Route C: Resolvent comparison via determinantal identities}

\textbf{Goal.}
Use the integrable structure of $\KAi$
(Its--Izergin--Korepin--Slavnov form,
as in Paper~13cc) to compute
$\det(I - z\DAi)$ explicitly for small $z > 0$,
and deduce positivity from the determinantal expansion.

\textbf{Mechanism.}
$\KAi$ is an IIKS operator with explicit vectors
$(\Ai(x), \Ai'(x))^\top$.
The perturbed Fredholm determinant
$\det(I - z\DAi) = \det(I - z\PAi A^{\mathrm{arith}}\PAi
+ z\KAi)$
can be expanded as a Fredholm series;
the first-order term is
$-z\,\mathrm{tr}(\DAi)$,
which is related to the integrated mismatch
$\int_0^R (A^{\mathrm{arith}}(x,x) - \KAi(x,x))\,dx$.
If this trace is positive, it gives evidence for
(but not a proof of) $\inf\sigma(\DAi) > 0$.

\textbf{Status.}
$\KAi(x,x) = \mathrm{tr}\,\KAi = \int_0^\infty K^{\mathrm{Airy}}(x,x)\,dx$
is explicitly computable via the Painlevé~II representation~\cite{TracyWidom1994}.
The term $A^{\mathrm{arith}}(x,x) = \lim_{N}\frac{1}{N}\sum_j |\Ai(p_j^{\mathrm{Ai}} - x)|^2$
requires estimates on the Airy function at prime arguments.

%%=================================================================
\section{Relation to First Airy Zero}
\label{sec:zero}
%%=================================================================

\begin{conjecture}[Connection to $a_1$]\label{conj:a1}
The spectral gap satisfies
\[
  c_2^{\mathrm{Airy}}
  \asymp |a_1|^{-1}
  \approx 0.428,
\]
where $a_1 \approx -2.338$ is the first zero of $\Ai$.
\end{conjecture}

\begin{remark}[Motivation]
The first Airy zero governs:
(i) the largest eigenvalue of $\KAi$ on $[a_1, \infty)$
(Tracy--Widom edge);
(ii) the oscillation period of $|\Ai(t)|^2$ near $t = 0$;
(iii) the spacing of the smallest prime-scaled sampling
points in the Airy window.
All three are controlled by $|a_1|$,
suggest a common scale for the spectral gap.
This conjecture is numerically checkable via the
Bornemann algorithm~\cite{Bornemann2010}.
\end{remark}

%%=================================================================
\section{Independence Certificate}
\label{sec:independence}
%%=================================================================

This problem is \textbf{independent of O4} (frame stability).
Specifically:
\begin{itemize}
  \item The statement of Problem~\ref{prob:O5} involves
    only the Airy operator $\KAi$, the arithmetic sampling form
    $A^{\mathrm{arith}}$ in the Airy scaling limit,
    and the Airy edge subspace.
  \item No Gram matrix, PSWF basis, edge block index
    $[(1-\delta)N, N)$, or Paley--Wiener frame condition appears.
  \item Routes~A--C above use Airy function analysis,
    Fredholm determinant theory, and IIKS structure;
    none uses sampling stability or prime gaps.
\end{itemize}

Resolution of O5 does not affect and does not require
resolution of O4.

%%=================================================================
\section{Open Sub-Problems}
\label{sec:subproblems}
%%=================================================================

\begin{problem}[Trace computation]\label{sub1}
Compute $\mathrm{tr}(\DAi)
= \int_0^R (A^{\mathrm{arith}}(x,x) - \KAi(x,x))\,dx$
numerically and establish its sign rigorously.
\end{problem}

\begin{problem}[$R$-independence]\label{sub2}
Show that $c_2^{\mathrm{Airy}}$ is bounded away from zero
uniformly in $R > 0$.
Equivalently: the spectral gap of $\DAi$ does not
collapse as the edge window grows.
\end{problem}

\begin{problem}[Verify Conjecture~\ref{conj:a1} numerically]%
\label{sub3}
Using the Bornemann algorithm~\cite{Bornemann2010},
compute $\inf\sigma(\DAi)$ numerically for $R = 5, 10, 20$
and compare with $|a_1|^{-1} \approx 0.428$.
\end{problem}

\begin{problem}[Density universality]\label{sub4}
Is $c_2^{\mathrm{Airy}}$ the same for all sampling
sequences $\{p_j^{\mathrm{Ai}}\}$ with the same Beurling density,
or does it depend on arithmetic properties of the primes?
(I.e., is the gap a universal constant or a number-theoretic invariant?)
\end{problem}

\begin{thebibliography}{99}
\bibitem{TracyWidom1994}
C.~Tracy and H.~Widom,
\textit{Level-spacing distributions and the Airy kernel},
Comm.\ Math.\ Phys.\ \textbf{159} (1994), 151--174.
\bibitem{Bornemann2010}
F.~Bornemann,
\textit{On the numerical evaluation of distributions in random
matrix theory: a review},
Markov Processes Relat.\ Fields \textbf{16} (2010), 803--866.
\bibitem{ItsIzergin1990}
A.~R.~Its, A.~G.~Izergin, V.~E.~Korepin, N.~A.~Slavnov,
\textit{Differential equations for quantum correlation functions},
Int.\ J.\ Mod.\ Phys.\ B \textbf{4} (1990), 1003--1037.
\bibitem{Deift1999}
P.~Deift, \textit{Orthogonal Polynomials and Random Matrices:
A Riemann--Hilbert Approach}, AMS, 1999.
\bibitem{PaperXXII}
S.~Schmalnauer,
\textit{Edge Coercivity and the Airy Boundary Operator}
(Paper~XXII), preprint, 2026.
\bibitem{Paper13cc}
S.~Schmalnauer,
\textit{Airy Universality for PSWF: Exact IIKS Structure
and RHP Reduction} (paper13cc), preprint, 2026.
\end{thebibliography}

\end{document}
